%-------------------------
% Resume in Latex
% Author : Tejas Khairnar
%------------------------

%---- Required Packages and Functions ----

\documentclass[a4paper,11pt]{article}
\usepackage{latexsym}
\usepackage{xcolor}
\usepackage{float}
\usepackage{ragged2e}
\usepackage[empty]{fullpage}
\usepackage{wrapfig}
\usepackage{lipsum}
\usepackage{tabularx}
\usepackage{titlesec}
\usepackage{geometry}
\usepackage{marvosym}
\usepackage{verbatim}
\usepackage{enumitem}
\usepackage[hidelinks]{hyperref}
\usepackage{fancyhdr}
\usepackage{multicol}
\usepackage{graphicx}
\usepackage{cfr-lm}
\usepackage[T1]{fontenc}
\usepackage{fontawesome}
\usepackage{orcidlink}
\colorlet{orcidlogocol}{black}
\usepackage[most]{tcolorbox}
\usepackage{makecell}
\usepackage{arydshln}


% Set the margins
\geometry{left=0.85cm, top=0.8cm, right=0.85cm, bottom=0.6cm}


% Sections formatting
\titleformat{\section}{
  \vspace{-4pt}\scshape\raggedright\large
}{}{0em}{}[\color{black}\titlerule \vspace{-7pt}]


%-------------------------
% Custom commands

\newcommand{\resumePOR}[3]{
\vspace{0.5mm}\item[]
    \begin{tabular*}{\textwidth}[t]{l@{\extracolsep{\fill}}r}
    \hspace{-3mm}{#1}:\hspace{1mm} & \hspace*{0pt}\hfill{\footnotesize{ #3}} \vspace{-0.5mm}\\ \hspace{-2.9mm}#2 
    \end{tabular*}
    \vspace{0mm}
}

\newcommand{\resumeExp}[4]{
\vspace{2mm}\item[]
    \begin{tabular*}{\textwidth}[t]{l@{\extracolsep{\fill}}r}
        \hspace{-4.4mm} \textbf{#1} \vspace{0mm} & {\footnotesize\emph{#3}}\vspace{0mm}\\
        \hspace{-4.3mm} \footnotesize{\text{#2}} \vspace{2mm} &\footnotesize{#4}
    \end{tabular*}
    \vspace{-6.1mm}
}


\newcommand{\resumeProject}[4]{
\vspace{3mm}\item[]
    \begin{tabular*}{\textwidth}[t]{l@{\extracolsep{\fill}}r}
        \hspace{-4.4mm} \small\textbf{#1} \vspace{0mm} & {\footnotesize{#3}}\vspace{-1mm}\\
        \hspace{-4.3mm} \footnotesize{\text{#2}}\vspace{2mm} & \footnotesize{#4} \vspace{0mm}
    \end{tabular*}
    \vspace{-6mm}
}

\newcommand{\resumeEdu}[4]{
\vspace{0mm}\item[]
    \begin{tabular*}{\textwidth}[t]{l@{\extracolsep{\fill}}r}
        \hspace{-4.3mm} \small\textbf{#1} & \footnotesize{#3}\vspace{-1mm} \\
        \hspace{-4.3mm} \footnotesize{#2} & \footnotesize{#4}
    \end{tabular*}
    \vspace{-3.2mm}
}

\newcommand{\resumeAchieve}[3]
{
\hspace{-3.1mm}\textbf{ #1} & {#2} & \hspace{3mm}\footnotesize{#3}
\vspace{0mm}\\
}


\renewcommand{\labelitemi}{$\vcenter{\hbox{\tiny$\bullet$}}$}

\newcommand{\resumeSubHeadingListStart}{\begin{itemize}[leftmargin=*,labelsep=0mm,itemsep=-2.5mm]}

\newcommand{\resumeItemListStart}{\begin{justify}\begin{itemize}[leftmargin=3ex, rightmargin=2ex, noitemsep,labelsep=1.2mm,itemsep=0mm]\small}
\newcommand{\resumePubItemListStart}{\begin{justify}\begin{itemize}[leftmargin=3ex, rightmargin=2ex, labelsep=1.2mm,itemsep=1.0mm]\small}
\newcommand{\resumePubHeading}[1]{\vspace{1mm}\item[] \textbf{#1}\vspace{-2.5mm}}

\newcommand{\resumeSubHeadingListEnd}{\end{itemize}\vspace{-2mm}}
\newcommand{\resumeItemListEnd}{\end{itemize}\end{justify}\vspace{-1.5mm}}
\newcommand{\resumePubItemListEnd}{\end{itemize}\end{justify}\vspace{-1mm}}

\renewcommand{\arraystretch}{1}

\newcolumntype{L}{>{\raggedright\arraybackslash}X}%
\newcolumntype{R}{>{\raggedleft\arraybackslash}X}%
\newcolumntype{C}{>{\centering\arraybackslash}X}%
%---- End of Packages and Functions ------



%-------------------------------------------
%%%%%%  CV STARTS HERE  %%%%%%%%%%%

\newcommand{\name}{Audrey Houghton (she/her)} % Your Name
\newcommand{\course}{} % Your Course
\newcommand{\roll}{xxxxxxxx} % Your Roll No.
\newcommand{\phone}{612-462-1481} % Your Phone Number
\newcommand{\emaila}{audreymhoughton@gmail.com} %Email 1
\newcommand{\github}{audreymhoughton} %Github
\newcommand{\website}{https://example.com} %Website
\newcommand{\linkedin}{audreyhoughton} %linkedin
\newcommand{\orcid}{0000-0001-9426-7969} %ORCID



\begin{document}
\fontfamily{cmr}\selectfont

%----------HEADING-----------------
\begin{center}
    \LARGE{\textbf{\name}}
\end{center}
\vspace{-5mm}
\begin{center}
    \small{\href{https://github.com/\github}{\faGithub \hspace{0.2mm} github.com/audreymhoughton} |  \href{https://www.linkedin.com/in/\linkedin/}{\faLinkedinSquare \hspace{0.2mm} linkedin.com/in/audreyhoughton} | \href{mailto:\emaila}{\faSend \hspace{0.2mm} \emaila} \\
    \faPhone \hspace{0.2mm} +1-\phone } | \href{https://scholar.google.com/citations?hl=en&user=ik59nO0AAAAJ&view}{\faGraduationCap Google Scholar} | \orcidlinkc{\orcid}
\end{center}
\vspace{-6mm}
\section{Research Interests}
\resumeItemListStart
\item[] Scientific software engineering, astronomical data processing, high-performance computing, machine learning for scientific data analysis, reproducible research workflows, large-scale data pipelines, and computational imaging. My work has contributed to more than 30 scientific publications spanning astronomy, neuroimaging, psychiatry, and scientific software.
\resumeItemListEnd
\vspace{-2mm}

\vspace{0mm}
\section{Education}

\resumeSubHeadingListStart
\resumeEdu
{University of Montana - Missoula, MT} 
{Bachelors of Arts in Physics} 
{Aug 2013 - May 2017} 
{\href{https://www.umt.edu/physics-astronomy/}{umt.edu/physics-astronomy}}
\resumeSubHeadingListEnd
\vspace{0mm}

\vspace{-2mm}

\section{Research Experience}
\resumeSubHeadingListStart

\resumeExp
{University of Minnesota – Department of Psychiatry \& Behavioral Sciences}
{Research Software Engineer III}
{Dec 2025 - Present}
{Minneapolis, MN (Remote)}
\resumeItemListStart
\item[$\bullet$] Develop and maintain Python analysis pipelines for GIS-linked electronic health record data in psychiatric research
\item[$\bullet$] Build geospatial feature-engineering workflows in Python/Jupyter for county- and tract-level overlays, geocoding quality control, and FIPS recovery from external APIs
\item[$\bullet$] Build publication-grade clinical prediction workflows with LightGBM, repeated stratified cross-validation, calibration selection, and bootstrap/permutation robustness checks
\item[$\bullet$] Lead statistical result interpretation and methods reporting for manuscripts and grant submissions across interdisciplinary teams
\resumeItemListEnd

\resumeExp
{Masonic Institute of the Developing Brain}
{Research Software Engineer II (promoted to III)}
{Jan 2021 - Apr 2024}
{Minneapolis, MN (Remote)}
\resumeItemListStart
\item[$\bullet$] Developed Python wrappers leveraging AWS S3 storage and SLURM optimization to accelerate image processing up to 12×.
\item[$\bullet$] Engineered a containerized infant MRI segmentation application achieving 600× faster processing and improved segmentation accuracy.
\item[$\bullet$] Built ETL pipelines for petabyte-scale neuroimaging datasets, emphasizing reproducibility and quality control.
\item[$\bullet$] Authored laboratory documentation and technical curriculum for neuroimaging software and HPC workflows.
\item[$\bullet$] Co-designed and instructed a three-month MRI Data Processing Workshop covering Bash, Linux, SLURM, BIDS, containers, and MRI analysis workflows.
\item[$\bullet$] Founded and led weekly Bash training sessions for research staff and students.
\item[$\bullet$] Founded and organized the interdisciplinary "Brainy Brains" seminar series to promote neuroscience learning across research disciplines.
\item[$\bullet$] Mentored two University of Minnesota UROP/URA students in MRI data processing and reproducible workflows using ABCD and HCP-YA datasets.
\item[$\bullet$] Processed and maintained workflows for major neuroimaging cohorts including ABCD, BCP, HCP-Development, HCP-Young Adult, ABIDE I, and ABIDE II.
\resumeItemListEnd

\resumeExp
{Rutgers University Astrophysics Department}
{Research Assistant}
{May 2016 -- Aug 2016}
{New Brunswick, NJ}
\resumeItemListStart
\item[$\bullet$] Refactored legacy computational astrophysics software from IDL to Python, improving maintainability and accessibility for future research.
\item[$\bullet$] Developed and optimized galactic collision simulations through parameter estimation and computational modeling.
\item[$\bullet$] Performed analysis and validation of simulation outputs to evaluate model performance and scientific accuracy.
\item[$\bullet$] Collaborated with faculty and researchers on computational astrophysics projects using scientific programming and numerical methods.
\item[$\bullet$] Presented research findings through an oral presentation to faculty and researchers in the Rutgers Department of Physics and Astronomy and a poster presentation at the interdisciplinary Summer Research Symposium Featuring REU and RiSE Scholars.
\resumeItemListEnd

\resumeExp
{The MINERVA Project}
{Research Assistant}
{Jan 2015 -- May 2017}
{Missoula, MT}
\resumeItemListStart
\item[$\bullet$] Developed a Python-based photometric data reduction pipeline for detecting transiting exoplanets from time-domain FITS imaging data.
\item[$\bullet$] Implemented CCD calibration workflows, including bias, dark, and flat-field correction, followed by aperture and differential photometry for high-precision light curve generation.
\item[$\bullet$] Utilized Astropy, Photutils, and the scientific Python ecosystem to automate astronomical image processing, quality control, and photometric analysis.
\item[$\bullet$] Supported exoplanet transit detection and photometric follow-up as part of the MINERVA exoplanet research program.
\item[$\bullet$] Collaborated with faculty and student researchers to develop and improve scientific software supporting astronomical research.
\item[$\bullet$] Presented pipeline development and research results at the 2016 Montana Space Grant Consortium Conference.
\item[$\bullet$] Developed software supporting published astronomical research, including a refereed journal article and an American Astronomical Society conference abstract.
\resumeItemListEnd

\resumeSubHeadingListEnd
\vspace{-2mm}

%-----------PROJECTS-----------------
\section{Industry Experience}
\resumeSubHeadingListStart

\resumeExp
{CoreWeave, Inc.}
{Fleet Operations Engineer}
{Apr 2024 - Sept 2024}
{Remote}
\resumeItemListStart
\item[$\bullet$] Provisioned and operated thousands of NVIDIA GPU nodes in Kubernetes using Ansible, AWX, Grafana, and Go, achieving 1-2 week onboarding and 6+ months average production uptime
\item[$\bullet$] Diagnosed and resolved hardware and software failures across thousands of GPU compute nodes supporting large-scale AI workloads.
\item[$\bullet$] Partnered with 10+ teams—including networking, hardware, operations, and support—to streamline troubleshooting, accelerate incident resolution, and document SOPs and post-mortems, improving efficiency and reliability
\item[$\bullet$] Identified points in the node lifecycle where nodes were lost, implementing solutions that saved the company millions in revenue
\resumeItemListEnd

\resumeExp
{Calyxt}
{Laboratory Assistant \& Data Specialist}
{Oct 2019 - Aug 2020}
{Roseville, MN}
\resumeItemListStart
\item[$\bullet$] Developed Python scripts for laboratory data tracking, quality control, and analysis reporting
\resumeItemListEnd

\resumeSubHeadingListEnd
\vspace{-2mm}

\section{Teaching \& Mentoring}
\resumeSubHeadingListStart

\resumeExp
{University of Minnesota Undergraduate Research Opportunities Program (UROP/URA)}
{Undergraduate Research Mentor}
{2023}
{Minneapolis, MN}
\resumeItemListStart
\item[$\bullet$] Mentored two undergraduate researchers in neuroimaging data processing.
\item[$\bullet$] Trained students to process Human Connectome Project Young Adult (HCP-YA) and ABCD MRI datasets.
\item[$\bullet$] Introduced reproducible workflows and quality-control procedures for large neuroimaging datasets.
\resumeItemListEnd

\resumeExp
{Co-Instructor, MRI Data Processing Workshop}
{Masonic Institute for the Developing Brain, University of Minnesota}
{2022}
{Minneapolis, MN}
\resumeItemListStart
\item[$\bullet$] Co-designed and taught a three-month workshop introducing students and researchers to neuroimaging data processing and high-performance computing.
\item[$\bullet$] Developed instructional materials covering Bash/Linux, SLURM, S3 storage, BIDS standards, containers, quality assurance, and reproducible MRI processing workflows.
\item[$\bullet$] Held weekly office hours and provided technical mentorship to participants from diverse academic backgrounds.
\item[$\bullet$] The workshop was later adapted into an official University of Minnesota credit-bearing course.
\resumeItemListEnd

\resumeExp
{Mathematics Teacher}
{Imagine Prep Surprise}
{Jul 2017 - Jul 2018}
{Surprise, AZ}
\resumeItemListStart
\item[$\bullet$] Taught Algebra II, Precalculus, Calculus, and Cambridge AS/A-Level Mathematics to high school students.
\item[$\bullet$] Developed lesson plans, assessments, and instructional materials.
\item[$\bullet$] Contributed to improving Algebra proficiency from approximately 11\% to 70\% during my tenure.
\resumeItemListEnd

\resumeExp
{Teaching Assistant - Observational Astronomy}
{University of Montana}
{Aug 2016 - Dec 2016}
{Missoula, MT}
\resumeItemListStart
\item[$\bullet$] Assisted instruction for the undergraduate Observational Astronomy course.
\item[$\bullet$] Supervised observing sessions and telescope operations.
\item[$\bullet$] Delivered a 45-minute lecture on Python-based astronomical data reduction.
\item[$\bullet$] Assisted students with programming and data analysis assignments.
\resumeItemListEnd

\resumeExp
{Mathematics Tutor}
{Huntington Learning Center}
{Jul 2019 - Oct 2019}
{Plymouth, MN}
\resumeItemListStart
\item[$\bullet$] Provided one-on-one mathematics tutoring for K-12 students.
\item[$\bullet$] Adapted instruction to students with varying backgrounds and learning styles.
\resumeItemListEnd

\resumeExp
{Mathematics Tutor}
{Tutor Doctor}
{Jan 2019 - May 2019}
{Loveland, CO}
\resumeItemListStart
\item[$\bullet$] Delivered individualized mathematics tutoring to middle and high school students.
\item[$\bullet$] Prepared students for coursework and standardized examinations.
\resumeItemListEnd

\resumeExp
{Planetarium Outreach Educator}
{University of Montana Physics Department}
{Aug 2015 - May 2017}
{Missoula, MT}
\resumeItemListStart
\item[$\bullet$] Presented public planetarium shows introducing astronomy and physics concepts.
\item[$\bullet$] Communicated scientific concepts to audiences of all ages.
\item[$\bullet$] Supported departmental outreach activities to promote STEM education.
\resumeItemListEnd

\resumeSubHeadingListEnd
\vspace{-2mm}

\section{Publications}
\resumeSubHeadingListStart

\resumePubHeading{Peer-Reviewed Journal Articles}
\resumePubItemListStart
\item[] Mehta, K., et al. (2024). XCP-D: A robust pipeline for the post-processing of fMRI data. \textit{Imaging Neuroscience}, 2, imag-2-00257.
\item[] Luo, A. C., et al. (2024). Functional connectivity development along the sensorimotor-association axis enhances the cortical hierarchy. \textit{Nature Communications}, 15(1), 3511.
\item[] Keller, A. S., et al. (2023). Personalized functional brain network topography is associated with individual differences in youth cognition. \textit{Nature Communications}, 14(1), 8411.
\item[] Pines, A., et al. (2023). Development of top-down cortical propagations in youth. \textit{Neuron}, 111(8), 1316-1330.
\item[] Keller, A. S., et al. (2024). A general exposome factor explains individual differences in functional brain network topography and cognition in youth. \textit{Developmental Cognitive Neuroscience}, 66, 101370.
\item[] Covitz, S., et al. (2022). Curation of BIDS (CuBIDS): A workflow and software package for streamlining reproducible curation of large BIDS datasets. \textit{NeuroImage}, 263, 119609.
\item[] Wilson, M. L., et al. (2019). First radial velocity results from the miniature exoplanet radial velocity array (MINERVA). \textit{Publications of the Astronomical Society of the Pacific}, 131(1005), 115001.
\item[] Zhao, C., et al. (2024). A reproducible and generalizable software workflow for analysis of large-scale neuroimaging data collections using BIDS apps. \textit{Imaging Neuroscience}, 2, imag-2-00074.
\item[] Byington, N., et al. (2023). Polyneuro risk scores capture widely distributed connectivity patterns of cognition. \textit{Developmental Cognitive Neuroscience}, 60, 101231.
\item[] Shafiei, G., et al. (2024). Generalizable links between borderline personality traits and functional connectivity. \textit{Biological Psychiatry}, 96(6), 486-494.
\item[] Holt-Gosselin, B., et al. (2024). Familial risk for depression moderates neural circuitry in healthy preadolescents to predict adolescent depression symptoms in the Adolescent Brain Cognitive Development (ABCD) study. \textit{Developmental Cognitive Neuroscience}, 68, 101400.
\item[] Feczko, E., et al. (2025). Baby Open Brains: An open-source dataset of infant brain segmentations. \textit{Scientific Data}, 12(1), 1423.
\item[] Giovinazzi, M. R., et al. (2020). The HD 217107 planetary system: Twenty years of radial velocity measurements. \textit{Astronomische Nachrichten}, 341(9), 870-878.
\resumePubItemListEnd

\resumePubHeading{Conference Proceedings \& Abstracts}
\resumePubItemListStart
\item[] Shang, Z., et al. (2022). Learning strategies for contrast-agnostic segmentation via SynthSeg for infant MRI data. In \textit{Proceedings of the International Conference on Medical Imaging with Deep Learning} (pp. 1075-1084).
\item[] Richie-Halford, A., et al. (2022). NiRV: The neuroimaging report viewer. In \textit{Proceedings of the Organization for Human Brain Mapping}.
\item[] Houghton, A., et al. (2017). Project MINERVA's follow-up on wide-field, small telescope photometry to identify exoplanets. \textit{American Astronomical Society Meeting Abstracts}, 229, 146.08.
\item[] Keller, A., et al. (2024). Understanding cognitive development by capturing complex, multidimensional, childhood environments and individual-specific patterns of functional brain network organization. \textit{Biological Psychiatry}, S85.
\item[] Keller, A., et al. (2023). Highlighting the role of multidimensional childhood environments in functional brain network organization and cognitive development. \textit{Neuropsychopharmacology}, 48, 134-135.
\item[] Bertolero, M., et al. (2022). Multivariate patterns of functional connectivity are linked to borderline-spectrum symptoms in young adulthood and youth. \textit{Biological Psychiatry}, 91(9), S170-S171.
\resumePubItemListEnd

\resumePubHeading{Preprints \& Manuscripts}
\resumePubItemListStart
\item[] Goncalves, M., et al. (2025). FMRIPrep Lifespan: Extending a robust pipeline for functional MRI preprocessing to developmental neuroimaging. \textit{bioRxiv}.
\item[] Feczko, E., et al. (2024). Baby Open Brains: An open-source repository of infant brain segmentations. \textit{bioRxiv}.
\item[] Coffman, C., et al. (2024). Heritability estimation of subcortical volumes in a multi-ethnic multi-site cohort study. \textit{bioRxiv}.
\item[] Shafiei, G., et al. (2023). Generalizable links between symptoms of borderline personality disorder and functional connectivity. \textit{bioRxiv}.
\item[] Randolph, A. C., et al. (2026). Mental health outcomes of foster and adopted individuals with adverse childhood experiences: A validation of known risks using EHR data. \textit{medRxiv}, 2026.05.28.26354276.
\item[] Caola, L., et al. (2026). Before birth, beyond childhood: Understanding the influence of prenatal substance exposure on psychiatric diagnoses. \textit{medRxiv}, 2026.05.27.26354275.
\item[] Mahmoudi, M., et al. (2026). Cortical thickness and curvature in autism and ADHD: A mega-analysis. \textit{bioRxiv}.
\resumePubItemListEnd

\resumeSubHeadingListEnd
\vspace{-2mm}

\section{Research Software}
\resumeSubHeadingListStart

\resumeProject
{Sales Summary Pipeline} 
{Independent Automation Project} 
{2025} 
{[Proprietary]}
\resumeItemListStart
\item[$\bullet$] Built a proof-of-concept pipeline integrating Salesforce logging, CSV rotation, and Microsoft Teams alerts with compliant web scraping (robots.txt, throttling, retries)
\resumeItemListEnd

\resumeProject
{CoreWeave Internal Documentation} 
{CoreWeave, Inc.} 
{2024} 
{[Proprietary]}
\resumeItemListStart
\item[$\bullet$] Initiated and built company-wide documentation from the ground up, replacing scattered one-off documents with an organized, centralized Notion knowledge base covering terminology, SOPs, and team workflows
\resumeItemListEnd

\resumeProject
{SLURM Wrappers} 
{High-Performance Computing S3 and Disk SLURM-Compliant Wrappers} 
{2021 - 2024} 
{\href{https://github.com/DCAN-Labs/SLURM_wrappers}{github.com/DCAN-Labs/SLURM\_wrappers}}
\resumeItemListStart
\item[$\bullet$] Bash/Python wrappers adaptable to any SLURM-based data pipeline
\resumeItemListEnd

\resumeProject
{ABCD Study Collection 3165 Data Processing} 
{NIH Data Archive (NDA)} 
{2021 - 2024} 
{\href{https://collection3165.readthedocs.io/en/stable/}{collection3165.readthedocs.io}}
\resumeItemListStart
\item[$\bullet$] Community dataset enabling access to petabytes of processed/unprocessed MRI data under BIDS standards
\resumeItemListEnd

\resumeProject
{DCAN Labs Informational Guide} %Project Name
{Minnesota Supercomputing Institute Data Engineering Documentation} %Project Name, Location Name
{2022 - 2024} %Event Dates
{\href{https://dcan-labs-informational-guide.readthedocs.io/en/latest/}{dcan-labs-informational-guide.readthedocs.io}} %Website
\resumeItemListStart
\item[$\bullet$] Comprehensive, continually-supported documentation of usage of the Minnesota Supercomputing Institute to process, analyze, and transfer MRI data
\resumeItemListEnd

\resumeProject
{CABINET} %Project Name
{Container-Linking Wrapper} %Project Name, Location Name
{2022 - 2023} %Event Dates
{\href{https://github.com/DCAN-Labs/CABINET}{github.com/DCAN-Labs/CABINET}} %Website
\resumeItemListStart
\item[$\bullet$] Python-based universal container-linking wrapper employing FAIR principles
\resumeItemListEnd

\resumeProject
{BIBSNet} 
{Infant MRI Brain Segmentation Application} 
{2022 - 2023} 
{\href{https://github.com/DCAN-Labs/BIBSnet}{github.com/DCAN-Labs/BIBSnet}}
\resumeItemListStart
\item[$\bullet$] Containerized infant MRI application; achieved 600x faster processing and +53\% accuracy
\resumeItemListEnd

\resumeSubHeadingListEnd

\section{Technical Skills}
\vspace{0.5mm}
\resumeItemListStart
\item[]\textbf{Programming:} Python, Bash, UNIX/Linux, Git/GitHub, LaTeX, Docker, Singularity
\item[]\textbf{Infrastructure \& Data:} AWS (S3), Ceph, SLURM, Kubernetes, SQL, ETL pipelines, batch processing, data validation, Globus, Grafana, Prometheus
\item[]\textbf{Statistics \& Modeling:} logistic regression, predictive modeling, LightGBM, calibration (Platt/isotonic), repeated stratified cross-validation, bootstrap/permutation inference, FDR correction (Benjamini-Hochberg), SHAP interpretability, ROC-AUC/PR-AUC/Brier metrics, causal inference \& discovery, GIS analysis, QGIS visualization, EHR integration, geocoding/FIPS recovery
\item[]\textbf{Scientific Imaging \& ML:} Neuroimaging/MRI segmentation, astronomical image processing (FITS, Astropy, Photutils), CCD calibration, differential photometry, time-domain imaging, deep learning inference workloads and compute planning (e.g., nnU-Net)
\item[]\textbf{Productivity \& Collaboration:} Jira, Linear, Kanban, Agile/Scrum, SOP documentation, Slack, Microsoft Teams, Google Workspace, Box, Salesforce API automation, compliant web scraping
\item[]\textbf{Additional Exposure:} Matlab, R, IDL, JavaScript/Node.js, Go, Ansible/AWX, RedHat, ReadTheDocs, Google Apps Script, GCP, Terraform, Datadog, CI/CD (GitHub Actions, CircleCI)
\resumeItemListEnd
\vspace{-3mm}

\section{Leadership \& Service}
\resumeItemListStart
\item[] \textbf{Founder \& Organizer, Brainy Brains Seminar Series} \\
\textit{Masonic Institute for the Developing Brain, University of Minnesota} | 2022--2024 \\
\hspace*{1.5em}$\bullet$ Founded and led a weekly interdisciplinary seminar where researchers from diverse backgrounds learned and presented topics in neuroscience. \\
\hspace*{1.5em}$\bullet$ Fostered cross-disciplinary collaboration and scientific communication across the laboratory.
\item[] \textbf{Founder \& Organizer, Weekly Bash Training Series} \\
\textit{Masonic Institute for the Developing Brain, University of Minnesota} | 2021--2024 \\
\hspace*{1.5em}$\bullet$ Founded and led weekly Bash and Linux training sessions for researchers and students. \\
\hspace*{1.5em}$\bullet$ Developed instructional materials covering command-line computing, scripting, and high-performance computing workflows. \\
\hspace*{1.5em}$\bullet$ Improved computational proficiency across the laboratory by providing hands-on technical instruction and ongoing support.
\item[] \textbf{Vice President, Society of Physics Students} \\
University of Montana | 2016--2017
\item[] \textbf{Founding Member, Women in Physics} \\
University of Montana
\item[] \textbf{Selected Participant, Rutgers University Undergraduate Summer Exchange Program}
\item[] \textbf{Recipient, University of Montana Global Leadership Initiative Award}
\resumeItemListEnd
\vspace{-2mm}

\section{Conference Presentations}
\resumeItemListStart
\item[] \textbf{Poster Presentation}, Montana Space Grant Consortium Conference, 2016. Presented research on the development of a Python pipeline for Project MINERVA.
\item[] \textbf{Poster Presentation}, Flux Congress, 2023. Presented CABINET, a container-linking wrapper for reproducible neuroimaging workflows.
\resumeItemListEnd
\vspace{-2mm}

\section{Professional Development}
\resumeItemListStart
\item[] \textbf{Leading on All Levels}, University of Minnesota (2021)
\item[] \textbf{Project Management Certificate Coursework}, University of Minnesota (2023)
\item[] \textbf{Data Management Series}, University of Minnesota (2023)
\resumeItemListEnd
\vspace{-2mm}

%-------------------------------------------
\end{document}

